\subsection*{General Objective}
Design and implement an ESP32-based IoT prototype to monitor, remotely control, and support 
the energy management of a domestic electric water heater. The system combines water-temperature, 
electrical-current, and water-flow measurements with an SSR, wireless communication, a dashboard, 
and a virtual-assistant interface.

\subsection*{Specific Objectives}
\begin{enumerate}[label=SO\arabic*.]
\item Acquire and calibrate water-temperature, electrical-current, and water-flow readings from the 
three sensors.

\item Process the readings on the ESP32 through temperature smoothing, RMS current calculation, and 
flow-pulse counting.

\item Transmit processed readings through Wi-Fi/MQTT and display the heater's status and measurements 
on a remotely accessible dashboard.

\item Test remote and automatic on/off switching through a suitably rated SSR, while keeping the heater's 
built-in protections in place.

\item Calculate estimated electrical energy, powered time, heating time, water-use time, water volume, 
and heating status from the measurements.

\item Implement configurable rules for a temperature target and prolonged powered operation without 
detected water use, with safeguards against unnecessary switching.

\item Implement voice or chat commands such as ``turn on the heater,'' 
``turn off the heater,'' and ``is the water hot?''

\item Validate sensor readings, dashboard information, and switching against independent measurements 
and the heater's observed operation.
\end{enumerate}

\subsection*{Scope}
\begin{itemize}
\item The project will deliver a functional IoT prototype focused on one domestic electric water heater, 
using the Rheem 50 L AI Smart model as the reference case. It will add an external monitoring and control 
layer to a heater that already has its own AI Smart functions and temperature display.

\item The ESP32 will read water temperature, alternating current, and water flow. Current will help 
distinguish a powered heater from one whose heating element is drawing power; flow will indicate water 
use at the chosen installation point.

\item The system will send processed measurements wirelessly and present temperature, flow, operating 
status, estimated energy consumption, and operating times in a dashboard.

\item Remote on/off control and local automatic rules will be tested using an SSR and an electrical 
assembly appropriate for the actual load. The temperature-target rule will use a measured temperature 
at a specified location, and a no-use timeout will be configurable for a storage heater.

\item The work includes calibration and experimental validation of readings, a protective enclosure 
for the prototype, and documentation of connections and system behavior.

\item An assistant interface for voice or chat commands and a local buzzer are optional additions. 
Work on the assistant may proceed in parallel, but integration depends on the three sensors, dashboard, 
and switching path functioning reliably.
\end{itemize}

\subsection*{Limitations}
\begin{itemize}
\item Validation will take place in a controlled test setting or pilot installation. The prototype 
is not a certified commercial product or a substitute for a qualified electrical installation.

\item The SCT-013 measures current, not active power directly. Energy values will be estimates 
based on nominal voltage and the assumed behavior of the heater, and will not have billing-grade accuracy.

\item A temperature probe on an accessible pipe or test reservoir may not measure the temperature 
throughout the 50 L tank. Its mounting point and the meaning of the displayed temperature will be 
documented.

\item The YF-S201 will initially be tested in a separate water circuit. Installation in a household 
water line depends on confirming the selected unit's temperature, pressure, fittings, and material 
suitability.

\item Control is limited to on/off switching; the project does not modify the heater's internal 
thermostat, thermal protections, AI Smart controller, or proportional heating power. External power 
switching must be checked against the heater's restart behavior.

\item The reference load is a single-phase 220 V electric heater. Three-phase installations and other 
household appliances are outside the proposed validation.

\item The virtual assistant is optional and depends on its chosen speech or chat service and internet 
connection. Local monitoring and protective rules should remain usable if the remote connection fails.

\item The project will use basic access control for remote commands, but a full cybersecurity assessment 
and long-term reliability certification are outside its scope.

\end{itemize}
